%---------------------------------------------------------------
% SLO: slovenski povzetek
% ENG: slovenian abstract
%---------------------------------------------------------------
\selectlanguage{slovene} % Preklopi na slovenski jezik
\phantomsection
\addcontentsline{toc}{chapter}{Povzetek}
\chapter*{Povzetek}

\noindent\textbf{Naslov:} \ttitle
\bigskip

Volumetrično sledenje potem je standardna tehnika za upodabljanje volumetričnih
podatkov, ki ohranja vse strukturne informacije v končni sliki, vendar je ta
pristop računsko zahtevnejši od upodabljanja površin. To je privedlo do
številnih optimizacijskih tehnik, od strategij medpomnjenja do nedavnih
pristopov s strojnim učenjem. Na podlagi predhodnega dela, ki je za napovedovanje
posredne sevalnosti prek preslikavanja fotonov uporabljalo nevronsko mrežo,
predlagamo metodo, ki preslikavanje fotonov nadomesti s sledenjem potem. S tem
odpravimo potrebo po ločeni fazi učenja in omogočimo sprotno učenje modela.
Naša metoda je vgrajena v obstoječi platformno agnostični volumetrični sledilnik
potem, ki deluje v brskalniku s pomočjo tehnologije \mbox{WebGPU}. Dosegamo
$2{,}06$-kratno skupno pohitritev v fazi računanja na grafičnem procesorju in
$1{,}28$-kratno izboljšavo časa sličice v primerjavi s sledenjem potem. Pokažemo,
da naša metoda ustvari slike z manj šuma in da so slike, pridobljene po
$0{,}5$ sekunde upodabljanja, bližje referenčnim slikam glede na zaznavno metriko
LPIPS kot slike, pridobljene s sledenjem potem.

\subsection*{Ključne besede}
\textit{\tkeywords}
\clearemptydoublepage

